\section*{Chapter \ref{ch:rs:transaction-costs-in-research} --- Transaction Costs in Research}

\begin{solution}{exo:rs:transaction-costs-in-research:1}
$2 + 10\,000 \times 0.7 \times 0.02 \times \sqrt{0.05} = 2 + 31.3 = 33.3$ basis points.
\end{solution}

\begin{solution}{exo:rs:transaction-costs-in-research:2}
$0.491/(5.33 \times 252) = 3.66$ basis points per unit traded; its half spread alone is 2, and impact at \$1 billion is far above the rest.
\end{solution}

\begin{solution}{exo:rs:transaction-costs-in-research:3}
Apart: $2q^{3/2}$; netted: $(2q)^{3/2} = 2^{3/2}q^{3/2}$; the netted trade costs $(2^{3/2} - 2)/2 = 41\%$ more.
\end{solution}

\begin{solution}{exo:rs:transaction-costs-in-research:4}
Its cost is dominated by the price's move while it works (63 basis points of standard deviation at 1\% participation against 14 of impact). The coefficient is an
average over orders: 500 orders pin $\eta$ to within about a third (0.24 around 1.09), 5\,000 to about 12\% (0.08 around 0.71), 50\,000 to 4\%.
\end{solution}

\begin{solution}{exo:rs:transaction-costs-in-research:5}
The optimiser trades whenever the forecast's change pays for the cost. At \$10 million costs are small, so it follows the forecast's daily noise, which is not
alpha, and pays 46.1\% a year to do so. At \$1 billion costs make it move slowly, which averages the noise; it even earns more before costs.
\end{solution}

\begin{solution}{exo:rs:transaction-costs-in-research:6}
For $v > 0$ the minimiser of $\frac12(x - v)^2 + s\lvert x\rvert + c\lvert x\rvert^{3/2}$ is non-negative; if $v \le s$ it is 0; otherwise the first-order condition
$x - v + s + \frac32 c\sqrt x = 0$ gives, with $y = \sqrt x$, $y^2 + \frac32 cy - (v - s) = 0$, so $y = \frac12(-\frac32 c + \sqrt{\frac94 c^2 + 4(v - s)})$ and
$x = y^2$; by symmetry the sign of $v$ is restored.
\end{solution}

\begin{solution}{exo:rs:transaction-costs-in-research:7}
Costs ignored: $-29.47$ raw, rising to 1.15 at fifty days and falling to 0.69 at two hundred. Costs inside: 1.80 raw, 2.43 at one day, 0.94 at two hundred. The
cost-aware optimiser already smooths the book through its reluctance to trade, so it needs only a little smoothing of the forecast; the naive book gets all its
smoothing from the forecast.
\end{solution}

\begin{solution}{exo:rs:transaction-costs-in-research:8}
The cost per unit traded grows as the square root of the size of each trade, so as a fraction of capital the costs rise by $\sqrt 2 = 1.41$ times when the fund
doubles; in dollars they rise $2^{3/2} = 2.83$ times, while the gross dollars only double. The fund's return after costs falls with size, and past some size its
dollar profit falls too (chapter~28).
\end{solution}

\begin{solution}{pb:rs:transaction-costs-in-research:1}
\begin{enumerate}
\item Half spread 2 basis points plus $0.7\sigma\sqrt{Q/V}$; at 1\% participation and 2\% volatility, 14 basis points of impact against a price noise of 63.
\item 500: $\eta = 1.09$ (0.24), exponent 0.59 (0.07); 5\,000: 0.71 (0.08), 0.48 (0.03); 50\,000: 0.67 (0.03), 0.49 (0.01).
\item Almgren and co-authors: a 3/5 power on about 700\,000 Citigroup orders; Tóth and co-authors: the square root.
\item Averages within bins remove most of the price noise before taking logarithms; the noise grows with the order's size, so a homoskedastic standard error is
wrong.
\item Turnover 5.33 a day, costs 1\,223\% a year, net Sharpe ratio $-29.47$.
\item $-0.31$, 0.06, 1.80 and 1.76.
\item It trades slowly and so averages the forecast's daily noise.
\item Costs make it trade 0.018 of capital a day; its Sharpe ratio before costs falls to 2.11 as the book cannot follow even the slow signals, and nets 1.76.
\item 8.27\% and 8.73\% of a team's capital a year apart; 16.04\% netted, a saving of 0.96 points (5.7\%).
\item Cost grows as the $3/2$ power of the trade: the combined trade is dearer than the two apart (2.94 points a year here).
\item Costs ignored: $-29.47$, $-19.18$, $-10.62$, $-2.57$, 0.10, 1.02, 1.15, 0.93, 0.69 for no smoothing and half-lives of 1, 2, 5, 10, 20, 50, 100, 200 days. Costs
inside: 1.80, 2.43, 2.37, 2.17, 1.94, 1.71, 1.40, 1.15, 0.94.
\item \textbf{Named result.} The net Sharpe ratio peaks at a half-life of fifty days (1.15) for the book that ignores costs and at one day (2.43) for the cost-aware
book, at \$1 billion.
\item The cost-aware book on the forecast smoothed over one day (2.43).
\item Small: its net ratio falls from 1.78 at \$10 million to 1.15 at \$1 billion and $-0.23$ at \$10 billion.
\item With the cost model, the fund size assumed, the costs a year, turnover, \omterm{def:rs:transaction-costs-in-research:smoothing}{break-even cost}, and net results at several sizes.
\item When trades follow each other within the impact's decay time, as the daily book's do: a propagator charges the second trade for the first's remaining
impact.
\item From its brokers' pre-trade models or published estimates, scaled by volume and volatility (invariance), and replace it with its own fits as fills come in.
\item Their capacities are not additive: trades in the same direction add up under a concave law, so the firm's capacity is less than the sum.
\item Both slow the book towards the persistent part of the forecast; the \omterm{def:rs:portfolio-construction-ii:aim}{aim portfolio} does it with known decay rates, smoothing with a chosen half-life.
\item In the model of the market, in the optimiser and in the signal, at the size the strategy will run.
\end{enumerate}
\end{solution}

\begin{solution}{iq:rs:transaction-costs-in-research:1}
Concavely: average impact grows roughly as the square root of the order's size relative to daily volume, scaled by volatility. Evidence: Almgren and co-authors on
about 700\,000 Citigroup orders (a 3/5 power for \omterm{def:rs:transaction-costs-in-research:impact}{temporary impact}), Tóth and co-authors (square root), and invariance studies of portfolio transitions.
\end{solution}

\begin{solution}{iq:rs:transaction-costs-in-research:2}
Measure its turnover and \omterm{def:rs:transaction-costs-in-research:smoothing}{break-even cost}, cost it at the fund size it will run with a fitted impact model, build it with costs inside the optimiser, smooth it,
and look at the decay of its \omterm{def:rs:anatomy-of-a-predictor:ic}{information coefficient}: a daily 2.5 before costs can be negative after them.
\end{solution}

\begin{solution}{iq:rs:transaction-costs-in-research:3}
Record each parent order's decision and \omterm{def:rs:simulation-versus-live:calibration}{arrival prices}, fills, size, participation, volatility and volume; regress cost (net of half spread) on
$\sigma\sqrt{Q/V}$ with robust errors and fit the exponent on binned averages; control for the market's move; use thousands of orders and re-fit regularly.
\end{solution}

\begin{solution}{iq:rs:transaction-costs-in-research:4}
As a second-order cone constraint ($t_i \ge \lvert d_i\rvert^{3/2}$ is representable with rotated cones), or, as here, by a proximal method: the cost is separable and
its proximal step has a closed form, a soft threshold followed by the root of a quadratic in $\sqrt{\lvert d\rvert}$.
\end{solution}

\begin{solution}{iq:rs:transaction-costs-in-research:5}
Net their trades internally, share the saving by standalone cost, and cost each strategy at its marginal cost to the firm: same-direction trades cost more
together than apart, so combined capacity is below the sum.
\end{solution}

\begin{solution}{iq:rs:transaction-costs-in-research:6}
If the gross return is $g$ a year on capital $A$ and costs as a fraction of capital grow as $k\sqrt A$ (trades scale with $A$, cost per unit with $\sqrt A$), the dollar
profit is $gA - kA^{3/2}$, maximised at $A^* = (2g/(3k))^2$, where the return on capital is $g/3$; beyond it every dollar added lowers the profit.
\end{solution}
